\documentclass[10pt,a4paper]{amsart}
\usepackage[margin=19mm]{geometry}
\usepackage{amsmath,amssymb,mathtools}
\usepackage[scr=rsfs,cal=euler]{mathalfa}
\usepackage{xfrac}
\usepackage[unicode,hidelinks,bookmarks=false]{hyperref}
\newcommand{\ie}{i.e.\ }
\newcommand{\cf}{cf.\ }
\newcommand{\resp}{respectively\ }
\newcommand{\Subsection}{\subsection}
\providecommand{\nobreakdash}{\nobreak\hbox{-}}
\setlength{\emergencystretch}{2em}
\multlinegap0pt
\usepackage{bm}
\usepackage{bbm}
\usepackage{dsfont}
\usepackage{xspace}
\usepackage{cancel}
\usepackage{mathtools}
\usepackage{amsthm}
\makeatletter
\@ifundefined{theorem}{\newtheorem{theorem}{Theorem}}{}
\@ifundefined{proposition}{\newtheorem{proposition}[theorem]{Proposition}}{}
\makeatother
\newcommand{\GalK}{\operatorname{Gal}(\overline{K}/K)}
\renewcommand{\div}{\mathrm{div}}
\title[Twists arising from torsion points]{Twists arising from torsion points}
\author{Lukas Novak}
\date{Source: arXiv:2510.08486v1 (CC BY 4.0)}
\begin{document}
\maketitle

\noindent\emph{Source and licence.} Twists arising from torsion points. Version arXiv:2510.08486v1 (CC BY 4.0), distributed under the Creative Commons Attribution 4.0 International licence, as recorded in the arXiv licence metadata for this exact version and shown on the abstract page https://arxiv.org/abs/2510.08486v1. Full rights notice: licenses/arxiv-2510.08486-CC-BY-4.0.txt.\par\smallskip

\noindent\textit{Editorial scope.} Counted unit: the Proposition whose source label is \texttt{prop\_zw\_def} in the article (source0.tex lines 486--488, source label \texttt{prop\_zw\_def}), delivered together with the article's own complete original proof of it (source0.tex lines 490--513, one \texttt{proof} environment of 24 lines). No mathematics is written, repaired or re-derived: statement and proof are the article's own text line for line. Required context retained uncounted: none. Every definition, display and label of those lines is retained unchanged. Omitted from this package and not redistributed: 1--485, 489--489, 514--684. Nothing inside a retained excerpt is omitted or altered: every retained body is delivered exactly as the article's lines are written, so no omission receipt of any kind accompanies this package. Citations: the retained excerpts carry no citation command at all, so no bibliography entry of the article is transcribed and no reference is redistributed. Reference adaptation: none was needed and none was made; every reference inside the delivered unit resolves inside it, so no unresolved reference or citation is produced. Printed slips kept literal: no printed word, symbol, hypothesis or number of the counted statement or of its proof was altered. Layout and class: the article's own class is \texttt{amsart} with packages including geometry, amssymb, amsthm, amsmath, graphicx, hyperref, multirow, array, todonotes, babel, lmodern, comment, bbold, mathtools; the wrapper is the lane's pinned \texttt{amsart} editorial wrapper at 10pt on A4, which restates the theorem-like environments the retained text needs and adds the article's own macro definitions that the retained text needs (\texttt{\textbackslash GalK}, \texttt{\textbackslash div}), with extra wrapper packages (bm, bbm, dsfont, xspace, cancel, mathtools). The delivered numbering is the unit's own. Assets: no retained span contains a figure environment, a graphics-inclusion command, a PDF-page inclusion, a table or a TikZ picture; no image file is copied into this package, the package declares no asset dependency and no image file is redistributed with it. Licence: version arXiv:2510.08486v1 of this article is distributed under the Creative Commons Attribution 4.0 International licence, as recorded in the arXiv licence metadata for this exact version and shown on the abstract page https://arxiv.org/abs/2510.08486v1. A full rights notice is delivered at licenses/arxiv-2510.08486-CC-BY-4.0.txt. Characters: every retained excerpt is pure ASCII, so the delivered engine, XeLaTeX with its default Latin Modern text font, reports no missing character.
\begin{proposition}\label{prop_zw_def}
    There exists functions $z$ and $w$ in the function field $K(H_S^d)$ with the above divisors.
\end{proposition}

\begin{proof}
    Since the divisor $(T)+(T+S)+\dotso+(T+(p-1)S)-(O)-(S)-\dotso -((p-1)S)$ has  degree $0$, its sum of points is $\displaystyle{\sum_{k=0}^{p-1}(T+kS) - \sum_{k=0}^{p-1}kS = O}$ and all the points in the divisor are $K$-rational, we know that there exists a function $f \in K(E)$ such that
    \[
        \div{(f)} = (T)+(T+S)+\dotso+(T+(p-1)S)-(O)-(S)-\dotso -((p-1)S).
    \]
    As functions $z$ and $f$ have the same divisor, it follows that $z=\lambda f$ for some constant $\lambda \in \overline{K}$.

    On the other hand, we also have that the divisor of the function $z\circ \tau_S$, where $\tau_S$ is the translation by the point $S$, is
    \[
        \div{(z\circ \tau_S)} = (T-S)+(T)+\dotso+(T+(p-2)S)-(-S)-(O)-\dotso -((p-2)S) = \div{(z)}
    \]
    Therefore, we have that there is a constant $c \in \overline{K}$ such that $z(P+S) = c z(P)$ for all points $P\in E$. 
    Since $S$ is a $p$-torsion point, we have $z(P) = z(P+pS)=c^p z(P)$ for all points $P \in E$. From there it follows that $c$ is a $p$-th root of unity. Thus, we have that $c=\zeta_p^k$ for some integer~$k$.

    The function $z$ will be in the function field $K(H_S^d)$ if and only if it is fixed by the twisted action action of $\GalK$ by the cocycle $\xi_d$. Let us take a $\sigma \in \GalK$ such that $\sigma(\sqrt[p]{d}) = \zeta_p \sqrt[p]{d}$. We want to choose the constant $\lambda$ such that $z^{\sigma^*}=z$. 
    
    We have that
    \[
        z^{\sigma^*}(P) = z^\sigma(P+S) = (\zeta_p^k z(P))^\sigma= \zeta_p^k \sigma(\lambda) f(P)
    \]
    and $z(P) = \lambda f(P)$ for all points $P \in E$. So, we have to choose the constant $\lambda$ such that $\zeta_p^k \sigma(\lambda) = \lambda$ holds. By choosing $\lambda = \dfrac{1}{\sqrt[p]{d^k}}$ one can easily check that the last equation is satisfied. With such previously chosen constant $\lambda$ one can also easily check that $z^{\sigma^*}=z$ for all the others $\sigma \in \GalK$ and therefore we have that $z \in K(H_S^d)$. 
    
    Using the same arguments as for the function $z$ we also get that there exists such a function $w \in K(H_S^d)$.
\end{proof}

\end{document}
